% !TeX program = xelatex
% UTF-8; compile twice with XeLaTeX. No external images or data files needed.
\documentclass[a4paper,landscape,10pt]{article}
\usepackage[margin=14mm,headheight=14pt,headsep=5mm,footskip=7mm]{geometry}
\usepackage{fontspec}
\setmainfont{DejaVu Sans}
\newfontfamily\ipafont{DejaVu Sans}
\usepackage[table]{xcolor}
\usepackage{array,longtable,booktabs,fancyhdr}
\definecolor{ink}{HTML}{18384D}
\definecolor{band}{HTML}{E7EFF4}
\pagestyle{fancy}\fancyhf{}
\fancyhead[L]{\small\color{ink}Parlez Vous QSO?! · Italiano / Italian}
\fancyhead[R]{\small Englische Phrase · Traduzione italiana · IPA}
\fancyfoot[L]{\scriptsize Vorlage: Italian v 1.0 · Credits: Giancarlo -- I7SWX}
\fancyfoot[R]{\small\thepage}
\setlength{\parindent}{0pt}
\setlength{\tabcolsep}{3mm}
\renewcommand{\arraystretch}{1.25}
\newcommand{\ipa}[1]{{\ipafont[#1]}}
\begin{document}
{\LARGE\bfseries\color{ink}Parlez Vous QSO?!}\hfill{\large Italiano / Italian}

\vspace{2mm}
{\small QSO-Sprachführer mit italienischer IPA-Lautschrift. Die Lautschrift bezieht sich auf die italienische Spalte.}

\vspace{2mm}
{\footnotesize \textbf{Leseschlüssel:} \ipa{ˈ} Hauptbetonung, \ipa{ˌ} Nebenbetonung;
\ipa{ɛ, ɔ} offene e-/o-Laute; \ipa{ɲ} wie „gn“ in „campagna“, \ipa{ʎ} wie „gli“ in „famiglia“.
Doppelte Konsonanten werden länger gesprochen. Schrägstriche trennen Alternativen; \ipa{…} markiert Platzhalter.}

\vspace{1mm}
\fontsize{9}{11.8}\selectfont
\begin{longtable}{@{}>{\raggedright\arraybackslash}p{13mm}>{\raggedright\arraybackslash}p{73mm}>{\raggedright\arraybackslash}p{75mm}>{\raggedright\arraybackslash}p{90mm}@{}}
\toprule
\textbf{Ref.} & \textbf{English phrase} & \textbf{Italiano} & \textbf{Lautschrift (IPA)}\\
\midrule
\endfirsthead
\toprule
\textbf{Ref.} & \textbf{English phrase} & \textbf{Italiano} & \textbf{Lautschrift (IPA)}\\
\midrule
\endhead
\midrule
\multicolumn{4}{r}{\scriptsize Fortsetzung auf der nächsten Seite}\\
\endfoot
\bottomrule
\endlastfoot
\rowcolor{band} 1.0 & \textbf{Opening Greetings} & \multicolumn{2}{l}{\textbf{Saluti Iniziali}} \\*
1.1 & I am very pleased to meet you. & Sono molto lieto di conoscerti. & \ipa{ˈsono ˈmolto ˈljɛto di koˈnoʃʃerti} \\
1.2 & Thank you for your call & Grazie per la tua chiamata & \ipa{ˈɡrattsje per la ˈtua kjaˈmata} \\
1.3 & Good morning & Buongiorno. & \ipa{bwɔnˈdʒorno} \\
1.31 & Good day & Buongiorno. & \ipa{bwɔnˈdʒorno} \\
1.32 & Good afternoon & Buon pomeriggio. & \ipa{bwɔn pomeˈriddʒo} \\
1.33 & Good evening & Buona sera. & \ipa{ˈbwɔna ˈsera} \\
1.4 & I know only a little Italian but I enjoy QSOs in Italian. & Parlo poco italiano ma mi piace fare i QSO in italiano. & \ipa{ˈparlo ˈpɔko itaˈljano ma mi ˈpjatʃe ˈfare i ku ɛsse ɔ in itaˈljano} \\
\rowcolor{band} 2.0 & \textbf{Operator} & \multicolumn{2}{l}{\textbf{Operatore}} \\*
2.1 & My name is Peter. & Il mio nome è Peter. & \ipa{il ˈmio ˈnome ɛ ˈpiːtə} \\
2.2 & I have been licensed since 1960 & Ho la licenza dal 1960. & \ipa{ɔ la liˈtʃɛntsa dal milleˌnɔveˈtʃɛnto sesˈsanta} \\
2.3 & I have been licensed for 43 years. & Sono 43 anni che ho la licenza. & \ipa{ˈsono kwaˌrantaˈtre ˈanni ke ɔ la liˈtʃɛntsa} \\
2.4 & I live in Sittingbourne. & Abito in Sittingbourne.\newline Il mio QTH è Sittingbourne. & \ipa{ˈabito in ˈsɪtɪŋbɔːn / il ˈmio ku ti akka ɛ ˈsɪtɪŋbɔːn} \\
2.5 & Sittingbourne is a village. & Sittingbourne è un paese. & \ipa{ˈsɪtɪŋbɔːn ɛ un paˈeze} \\
2.51 & Sittingbourne is a town. & Sittingbourne è una cittadina. & \ipa{ˈsɪtɪŋbɔːn ɛ ˈuna tʃittaˈdina} \\
2.52 & Sittingbourne is a city. & Sittingbourne è una città. & \ipa{ˈsɪtɪŋbɔːn ɛ ˈuna tʃitˈta} \\
2.6 & Sittingbourne is between London and Dover. & Sittingbourne si trova tra Londra e Dover. & \ipa{ˈsɪtɪŋbɔːn si ˈtrɔva tra ˈlondra e ˈdəʊvə} \\
2.7 & I live in the South of England. & Abito nel Sud dell’Inghilterra. & \ipa{ˈabito nel sud dellinɡilˈtɛrra} \\
2.71 & I live in the North of England. & Abito nel nord dell’Inghilterra. & \ipa{ˈabito nel nɔrd dellinɡilˈtɛrra} \\
2.72 & I live in the East of England. & Abito nell’est dell’Inghilterra. & \ipa{ˈabito nell ɛst dellinɡilˈtɛrra} \\
2.73 & I live in West of England. & Abito nell’ovest dell’Inghilterra. & \ipa{ˈabito nell ˈɔvest dellinɡilˈtɛrra} \\
2.74 & I live in South East England. & Abito nel sud est dell’Inghilterra. & \ipa{ˈabito nel sud ɛst dellinɡilˈtɛrra} \\
2.75 & I live in South West England. & Abito nel sud ovest dell’Inghilterra. & \ipa{ˈabito nel sud ˈɔvest dellinɡilˈtɛrra} \\
2.76 & I live in North East England. & Abito nel nord est dell’Inghilterra. & \ipa{ˈabito nel nɔrd ɛst dellinɡilˈtɛrra} \\
2.77 & I live in North West England. & Abito nel nord ovest dell’Inghilterra. & \ipa{ˈabito nel nɔrd ˈɔvest dellinɡilˈtɛrra} \\
2.78 & I live in central England. & Abito nel centro dell’Inghilterra. & \ipa{ˈabito nel ˈtʃɛntro dellinɡilˈtɛrra} \\
2.8 & I live in the country. & Abito in campagna. & \ipa{ˈabito in kamˈpaɲɲa} \\
2.81 & I live by the sea. & Abito vicino al mare. & \ipa{ˈabito viˈtʃino al ˈmare} \\
2.82 & I live near mountains. & Abito in montagna. & \ipa{ˈabito in monˈtaɲɲa} \\
2.9 & Sittingbourne is 50Km from London. & Sittingbourne è a 50Km da Londra. & \ipa{ˈsɪtɪŋbɔːn ɛ a tʃinˈkwanta kiˈlɔmetri da ˈlondra} \\
2.11 & I am retired. & Sono un pensionato. & \ipa{ˈsono un pensjoˈnato} \\
2.12 & My home call sign is M0ECS. & Il mio nominativo originario è M0ECS. & \ipa{il ˈmio nominaˈtivo oridʒiˈnarjo ɛ ɛmme ˈdzɛro e tʃi ɛsse} \\
\rowcolor{band} 3.0 & \textbf{Station} & \multicolumn{2}{l}{\textbf{Stazione}} \\*
3.1 & My radio is an Icom 706. & Il mio apparato è un Icom 706.\newline La mia stazione è un Icom 706. & \ipa{il ˈmio appaˈrato ɛ un ˈaɪkɒm ˌsɛtteˌtʃɛntoˈsɛi / la ˈmia statˈtsjone ɛ un ˈaɪkɒm ˌsɛtteˌtʃɛntoˈsɛi} \\
3.2 & My station is all homebrew (homemade). & La mia stazione è tutta autocostruita. & \ipa{la ˈmia statˈtsjone ɛ ˈtutta autokoˈstruita} \\
3.3 & I am running 5 Watts. & La mia potenza è di 5 Watt. & \ipa{la ˈmia poˈtɛntsa ɛ di ˈtʃiŋkwe vat} \\
3.4 & My antenna is a long wire. & La mia antenna è una longwire. & \ipa{la ˈmia anˈtɛnna ɛ ˈuna ˈlɒŋwaɪə} \\
3.41 & My antenna is a dipole. & La mia antenna è un dipolo. & \ipa{la ˈmia anˈtɛnna ɛ un diˈpɔlo} \\
3.42 & My antenna is a G5RV. & La mia antenna è una G5RV. & \ipa{la ˈmia anˈtɛnna ɛ ˈuna dʒi ˈtʃiŋkwe ɛrre vi} \\
3.43 & My antenna is a doublet. & La mia antenna è una doublet. & \ipa{la ˈmia anˈtɛnna ɛ ˈuna ˈdʌblət} \\
3.44 & My antenna is a delta loop. & La mia antenna è una delta loop. & \ipa{la ˈmia anˈtɛnna ɛ ˈuna ˈdɛlta luːp} \\
3.45 & My antenna is a vertical. & La mia antenna è una verticale. & \ipa{la ˈmia anˈtɛnna ɛ ˈuna vertiˈkale} \\
3.46 & My antenna is a log periodic. & La mia antenna è una log periodica. & \ipa{la ˈmia anˈtɛnna ɛ ˈuna lɔɡ peˈrjɔdika} \\
3.47 & My antenna is a Yagi beam. & La mia antenna è una yagi. & \ipa{la ˈmia anˈtɛnna ɛ ˈuna ˈjaɡi} \\
3.48 & My antenna is a 5 element beam. & La mia antenna è una 5 elementi. & \ipa{la ˈmia anˈtɛnna ɛ ˈuna ˈtʃiŋkwe eleˈmenti} \\
\rowcolor{band} 4.0 & \textbf{Conditions} & \multicolumn{2}{l}{\textbf{Condizioni}} \\*
4.1 & Thank you for my 599 report. & Grazie per il mio rapporto di 599. & \ipa{ˈɡrattsje per il ˈmio rapˈpɔrto di ˈtʃiŋkwe ˈnɔve ˈnɔve} \\
4.2 & The band conditions are good. & Le condizioni di banda sono buone. & \ipa{le konditˈtsjoni di ˈbanda ˈsono ˈbwɔne} \\
4.3 & The band conditions are bad. & Le condizioni di banda sono pessime. & \ipa{le konditˈtsjoni di ˈbanda ˈsono ˈpɛssime} \\
4.4 & There is strong interference here. & Ci sono delle forti interferenze. & \ipa{tʃi ˈsono ˈdelle ˈfɔrti interferˈɛntse} \\
4.5 & Can you change to the 80m band? & Puoi cambiare banda e passare su 80 metri? & \ipa{pwɔi kamˈbjare ˈbanda e pasˈsare su otˈtanta ˈmɛtri} \\
4.6 & Please change frequency to 3.558 MHz. & Puoi cambiare la frequenza su 3,558 MHz? & \ipa{pwɔi kamˈbjare la freˈkwɛntsa su tre ˈvirɡola ˈtʃiŋkwe ˈtʃiŋkwe ˈɔtto ˈmɛɡaɛrts} \\
\rowcolor{band} 5.0 & \textbf{General Queries and Comments} & \multicolumn{2}{l}{\textbf{Informazioni e Commenti Generici}} \\*
5.1 & What does <\#\#\#\#\#> mean please? & Per cortesia che cosa vuole dire <\#\#\#\#>? & \ipa{per korteˈzia ke ˈkɔza ˈvwɔle ˈdire …} \\
5.2 & Please wait a moment, I need to check my Italian dictionary. & Puoi attendere un momento che controllo il mio vocabolario d’italiano. & \ipa{pwɔi atˈtɛndere un moˈmento ke konˈtrɔllo il ˈmio vokaboˈlarjo ditaˈljano} \\
5.3 & How do you say, <\#\#\#\#\#> in Italian ? & Come si dice <\#\#\#\#\#> in italiano? & \ipa{ˈkome si ˈditʃe … in itaˈljano} \\
5.5 & My favourite modes of operation are CW / RTTY / PSK / MT63 / Olivia / Hellschreiber / Digital Modes / and SSB & I miei modi operativi preferiti sono CW / RTTY / PSK / MT63 / Olivia / Hellschreiber / Modi Digitali ed SSB. & \ipa{i ˈmjɛi ˈmɔdi operaˈtivi prefeˈriti ˈsono tʃi ˈdɔppja vu / ɛrre ti ti ˈipsilon / pi ɛsse kappa / ɛmme ti sesˌsantaˈtre / oˈlivja / ˈhɛlˌʃʁaɪbɐ / ˈmɔdi didʒiˈtali ed ɛsse ɛsse bi} \\
5.6 & My best DX is Australia. & Il mio migliore DX è l’Australia. & \ipa{il ˈmio miʎˈʎore di iks ɛ lauˈstralja} \\
5.7 & I enjoy working contests. & Mi piace partecipare ai contest. & \ipa{mi ˈpjatʃe partetʃiˈpare ai ˈkɒntɛst} \\
5.8 & I like working QRP. & Mi piace operare in QRP. & \ipa{mi ˈpjatʃe opeˈrare in ku ɛrre pi} \\
\rowcolor{band} 6.0 & \textbf{The Weather} & \multicolumn{2}{l}{\textbf{Il Tempo / il WX}} \\*
6.1 & It is sunny. & C’è il sole. & \ipa{tʃɛ il ˈsole} \\
6.2 & It is windy. & C’è vento. & \ipa{tʃɛ ˈvɛnto} \\
6.3 & It is raining. & Sta piovendo. & \ipa{sta pjoˈvɛndo} \\
6.4 & It is snowing. & Sta nevicando. & \ipa{sta neviˈkando} \\
6.5 & It is very hot. & Fa molto caldo. & \ipa{fa ˈmolto ˈkaldo} \\
6.6 & It is very cold. & Fa molto freddo. & \ipa{fa ˈmolto ˈfreddo} \\
6.7 & The temperature is 8C. & La temperatura è di 8 gradi. & \ipa{la temperaˈtura ɛ di ˈɔtto ˈɡradi} \\
6.8 & It is overcast. & È coperto.\newline È nuvoloso. & \ipa{ɛ koˈpɛrto / ɛ nuvoˈlozo} \\
6.9 & It is foggy. & C’è la nebbia. & \ipa{tʃɛ la ˈnɛbbja} \\
6.11 & It has been raining all day. & Ha piovuto tutto il giorno. & \ipa{a pjoˈvuto ˈtutto il ˈdʒorno} \\
6.12 & There is 20cm of snow here. & Ci sono 20cm di neve. & \ipa{tʃi ˈsono ˈventi tʃenˈtimetri di ˈneve} \\
\rowcolor{band} 8.0 & \textbf{QSL} & \multicolumn{2}{l}{\textbf{QSL}} \\*
8.1 & My QSL card is sure via the bureau. & Ti invierò la mia QSL via bureau (associazione). & \ipa{ti invjeˈrɔ la ˈmia ku ɛsse ɛlle ˈvia byˈʁo (assotʃatˈtsjone)} \\
8.2 & QSL via eQSL & QSL via e-QSL. & \ipa{ku ɛsse ɛlle ˈvia i ku ɛsse ɛlle} \\
8.3 & No eQSL please & No e-QSL, mi dispiace. & \ipa{nɔ i ku ɛsse ɛlle mi diˈspjatʃe} \\
8.4 & I will send you my card direct. & Ti invierò la QSL via diretta. & \ipa{ti invjeˈrɔ la ku ɛsse ɛlle ˈvia diˈrɛtta} \\
8.5 & What is your address? & Quale è il tuo indirizzo? & \ipa{ˈkwale ɛ il ˈtuo indiˈrittso} \\
8.6 & Please send your QSL card direct. & Per cortesia inviami la tua QSL via diretta. & \ipa{per korteˈzia ˈinvjami la ˈtua ku ɛsse ɛlle ˈvia diˈrɛtta} \\
8.7 & My address is... & Il mio indirizzo è …. & \ipa{il ˈmio indiˈrittso ɛ …} \\
8.8 & My address is shown on QSL.com / EHam.net / ….. & Il mio indirizzo è OK su qsl.com / EHam.net / ..... & \ipa{il ˈmio indiˈrittso ɛ ˌoˈkɛi su ku ɛsse ɛlle ˈpunto kɔm / i hæm ˈpunto nɛt / …} \\
8.9 & Sorry I do not send QSL cards. & Mi dispiace ma non spedisco QSL. & \ipa{mi diˈspjatʃe ma non speˈdisko ku ɛsse ɛlle} \\
8.11 & Please send QSL card via G1JQH. & Ti prego di inviarmi la QSL via G1JQH. & \ipa{ti ˈprɛɡo di inˈvjarmi la ku ɛsse ɛlle ˈvia dʒi ˈuno ˈi ˈluŋɡa ku akka} \\
8.12 & My WACRAL / GQRP / DXCC / FISTS number is 1234. & Il mio numero WACRAL / GQRP / DXCC / FISTS è 1234. & \ipa{il ˈmio ˈnumero ˈdɔppja vu a tʃi ɛrre ɛlle / dʒi ku ɛrre pi / di iks tʃi tʃi / fɪsts ɛ ˌmilledueˌtʃentoˌtrentaˈkwattro} \\
\rowcolor{band} 9.0 & \textbf{Final Greetings} & \multicolumn{2}{l}{\textbf{Saluti Finali}} \\*
9.1 & Thank you. & Grazie. & \ipa{ˈɡrattsje} \\
9.2 & Thank you for a very enjoyable contact. & Grazie per un piacevole contatto.\newline Grazie per un piacevole QSO. & \ipa{ˈɡrattsje per un pjaˈtʃevole konˈtatto / ˈɡrattsje per un pjaˈtʃevole ku ɛsse ɔ} \\
9.3 & I enjoyed our chat. & Mi ha fatto piacere chiacchierare con te. & \ipa{mi a ˈfatto pjaˈtʃere kjakkjeˈrare kon te} \\
9.4 & I enjoyed our QSO in Italian. & Mi è piaciuto il nostro QSO in italiano. & \ipa{mi ɛ pjaˈtʃuto il ˈnɔstro ku ɛsse ɔ in itaˈljano} \\
9.5 & It has been great fun. & Mi sono divertito molto. & \ipa{mi ˈsono diverˈtito ˈmolto} \\
9.6 & If you would like a regular contact with me, you can contact me via email. & Se vuoi avere dei collegamenti regolari con me mi puoi contattare via e-mail. & \ipa{se vwɔi aˈvere dei kolleɡaˈmenti reɡoˈlari kon me mi pwɔi kontatˈtare ˈvia ˈiːmeɪl} \\
9.7 & My email address is... & Il mio indirizzo e-mail è... & \ipa{il ˈmio indiˈrittso ˈiːmeɪl ɛ …} \\
9.8 & My email address is shown on QSL.com / EHam.net & Il mio indirizzo e-mail è riportato su qsl.com / EHam.net. & \ipa{il ˈmio indiˈrittso ˈiːmeɪl ɛ riporˈtato su ku ɛsse ɛlle ˈpunto kɔm / i hæm ˈpunto nɛt} \\
9.9 & I look forward to meeting you again. & Spero di incontrarti di nuovo. & \ipa{ˈspɛro di iŋkonˈtrarti di ˈnwɔvo} \\
9.11 & My best wishes to you and your family. & I miei migliori auguri a te ed alla tua famiglia. & \ipa{i ˈmjɛi miʎˈʎori auˈɡuri a te ed ˈalla ˈtua faˈmiʎʎa} \\
9.12 & Best wishes to your wife. & I migliori auguri a tua moglie. & \ipa{i miʎˈʎori auˈɡuri a ˈtua ˈmoʎʎe} \\
9.121 & Best wishes to your husband. & I migliori auguri a tuo marito. & \ipa{i miʎˈʎori auˈɡuri a ˈtuo maˈrito} \\
9.13 & Good luck. & Buona fortuna. & \ipa{ˈbwɔna forˈtuna} \\
9.14 & Good health. & Buona Salute. & \ipa{ˈbwɔna saˈlute} \\
9.15 & Goodbye & Arrivederci.\newline Ciao. & \ipa{arriveˈdertʃi / ˈtʃao} \\
9.16 & I wish you an enjoyable afternoon. & Ti auguro un buon pomeriggio.\newline Ti auguro una buona giornata. & \ipa{ti auˈɡuro un bwɔn pomeˈriddʒo / ti auˈɡuro ˈuna ˈbwɔna dʒorˈnata} \\
9.161 & I wish you an enjoyable evening. & Ti auguro una buona sera. & \ipa{ti auˈɡuro ˈuna ˈbwɔna ˈsera} \\
9.162 & I wish you an enjoyable weekend. & Ti auguro un buon fine settimana. & \ipa{ti auˈɡuro un bwɔn ˈfine settiˈmana} \\
9.17 & Enjoy your holiday. & Buone vacanze. & \ipa{ˈbwɔne vaˈkantse} \\
9.18 & Happy New Year. & Buon Anno. & \ipa{bwɔn ˈanno} \\
\rowcolor{band} 10.0 & \textbf{Country} & \multicolumn{2}{l}{\textbf{Nazione}} \\*
10.1 & Albania & Albania & \ipa{alˈbanja} \\
10.2 & Austria & Austria & \ipa{ˈaustrja} \\
10.3 & Belgium & Belgio & \ipa{ˈbɛldʒo} \\
10.4 & Belorussia & Bielorussia & \ipa{bjeloˈrussja} \\
10.5 & Bosnia \& Herzegovina & Bosnia - Herzegovina & \ipa{ˈbɔznja ertseˈɡɔvina} \\
10.6 & Bulgaria & Bulgaria & \ipa{bulˈɡarja} \\
10.7 & Croatia & Croazia & \ipa{kroˈattsja} \\
10.8 & Czech Republic & Repubblica Ceca & \ipa{reˈpubblika ˈtʃɛka} \\
10.9 & Denmark & Danimarca & \ipa{daniˈmarka} \\
10.11 & Eire & Irlanda & \ipa{irˈlanda} \\
10.12 & England & Inghilterra & \ipa{inɡilˈtɛrra} \\
10.13 & Estonia & Estonia & \ipa{eˈstɔnja} \\
10.14 & Finland & Finlandia & \ipa{finˈlandja} \\
10.15 & France & Francia & \ipa{ˈfrantʃa} \\
10.16 & Germany & Germania & \ipa{dʒerˈmanja} \\
10.17 & Gibraltar & Gibilterra & \ipa{dʒibilˈtɛrra} \\
10.18 & Greece & Grecia & \ipa{ˈɡrɛtʃa} \\
10.19 & Holland & Olanda & \ipa{oˈlanda} \\
10.21 & Hungary & Ungheria & \ipa{uŋɡeˈria} \\
10.22 & Iceland & Islanda & \ipa{iˈzlanda} \\
10.23 & Northern Ireland & Irlanda del Nord & \ipa{irˈlanda del nɔrd} \\
10.24 & Isle of Man & Isola di Man & \ipa{ˈizola di man} \\
10.25 & Italy & Italia & \ipa{iˈtalja} \\
10.26 & Latvia & Lettonia & \ipa{letˈtɔnja} \\
10.27 & Liechtenstein & Liechtenstein & \ipa{ˈliçtn̩ʃtaɪn} \\
10.28 & Lithuania & Lituania & \ipa{liˈtwanja} \\
10.29 & Luxembourg & Lussemburgo & \ipa{lussemˈburɡo} \\
10.31 & Macedonia & Macedonia & \ipa{matʃeˈdɔnja} \\
10.32 & Moldavia & Moldavia & \ipa{molˈdavja} \\
10.33 & Netherlands & Paesi Bassi & \ipa{paˈezi ˈbassi} \\
10.34 & Norway & Norvegia & \ipa{norˈvɛdʒa} \\
10.35 & Poland & Polonia & \ipa{poˈlɔnja} \\
10.36 & Portugal & Portogallo & \ipa{portoˈɡallo} \\
10.37 & Romania & Romania & \ipa{romaˈnia} \\
10.38 & Russian Federation & Russia & \ipa{ˈrussja} \\
10.39 & Scotland & Scozia & \ipa{ˈskɔttsja} \\
10.41 & Slovakia & Slovacchia & \ipa{zloˈvakkja} \\
10.42 & Slovenia & Slovenia & \ipa{zloˈvɛnja} \\
10.43 & Spain & Spagna & \ipa{ˈspaɲɲa} \\
10.44 & Svalbard & Isole Svalbard & \ipa{ˈizole ˈsvalbard} \\
10.45 & Sweden & Svezia & \ipa{ˈzvɛttsja} \\
10.46 & Switzerland & Svizzera & \ipa{ˈzvittsera} \\
10.47 & Turkey & Turchia & \ipa{turˈkia} \\
10.48 & Ukraine & Ucraina & \ipa{uˈkraina} \\
10.49 & Wales & Galles & \ipa{ˈɡalles} \\
\end{longtable}
\normalsize
\section*{Hinweise zur Übertragung und Aussprache}
\small\raggedright
\textbf{Vorlage und Umfang.} Grundlage ist „Italian v 1.0.doc“, mit dem Credit Giancarlo -- I7SWX.
Der Titel innerhalb der Quelle nennt abweichend „Italian v 21.0“; diese Versionsabweichung wird hier
festgehalten. Alle 142 Phrasen und Ländernamen einschließlich aller alternativen Formulierungen sind enthalten.
Die ursprüngliche Gliederung mit ihren Nummernsprüngen und dem fehlenden Abschnitt 7 bleibt erhalten.

\medskip
\textbf{Dokumentierte Korrekturen.} Apostrophe als Ersatz für Akzente wurden in der italienischen Spalte
berichtigt: „e’ / E’“ zu „è / È“, „citta’“ zu „città“ und „inviero’“ zu „invierò“.
„Sta’“ wurde in 6.3 und 6.4 als „Sta“ geschrieben. „modi operative“ wurde zu „modi operativi“ (5.5),
„temepratura“ zu „temperatura“ (6.7). Leerzeichen nach „nell’“ in 2.72 und 2.73 wurden entfernt;
mehrfache Leerzeichen vereinheitlicht. Das kaufmännische Und vor alternativen Formulierungen wurde
jeweils durch einen Zeilenumbruch ersetzt.

\medskip
\textbf{Inhaltliche Eigenheiten.} Die Übersetzungen bleiben ansonsten wie in der Quelle. So stehen in 2.82
„near mountains“ und „in montagna“, in 6.7 „8C“ und „8 gradi“ sowie in 8.3 „please“ und „mi dispiace“.
Die italienische Spalte verwendet überwiegend die vertraute Anrede „tu“. Die männlichen Formen
„lieto“, „pensionato“ und „divertito“ bleiben erhalten; für eine Sprecherin wären „lieta“, „pensionata“ und
„divertita“ zu verwenden. Historische Ländernamen und Webadressen wurden nicht modernisiert.
„Bosnia - Herzegovina“ bleibt als Schreibweise der Quelle erhalten; die Lautschrift verwendet die
italienische Aussprache von „Bosnia-Erzegovina“.

\medskip
\textbf{Transkriptionskonvention.} Die IPA-Spalte ist eine breite, manuell erstellte Aussprachehilfe für
Standarditalienisch, keine enge Transkription einer Aufnahme. \ipa{ˈ} steht vor der betonten Silbe;
bei Zahlwörtern und Zusammensetzungen kann \ipa{ˌ} eine Nebenbetonung markieren. Lange Konsonanten
werden durch Verdoppelung dargestellt, auch bei \ipa{ɲɲ} und \ipa{ʎʎ}; die Affrikaten \ipa{ts} und
\ipa{tʃ} sind ohne zusätzliche Bindebögen geschrieben. Die vorhersagbare Vokallänge des Italienischen
ist nicht eigens markiert. Wortübergreifende Konsonantenverdoppelung (raddoppiamento fonosintattico)
und regionale Varianten werden nicht systematisch angegeben.

\medskip
\textbf{Namen und Fachwörter.} Peter, Sittingbourne und Dover sowie Longwire, Doublet, Contest und
E-Mail sind mit englischen Ausspracheformen wiedergegeben; Hellschreiber und Liechtenstein mit deutschen.
Für Icom ist eine englische Aussprachevariante gewählt. Sonst werden italienische beziehungsweise
italienisch angepasste Formen verwendet; für „bureau“ steht eine französische Aussprache.
Bei Marken, Orten und Fachwörtern sind weitere gebräuchliche Aussprachen möglich.

\medskip
\textbf{Zahlen, Kürzel und Varianten.} Zahlen werden ausgeschrieben; 706 wird als „settecentosei“ gelesen,
599 als „cinque nove nove“ und 3,558 MHz als „tre virgola cinque cinque otto megahertz“.
Die Lautschrift folgt der italienischen Spalte. Rufzeichen und die meisten Kürzel werden mit italienischen
Buchstabennamen gelesen, nicht im internationalen Buchstabieralphabet: J als „i lunga“, W als „doppia vu“,
Y als „ipsilon“, V als „vi“. FISTS wird als Wort gesprochen, das e in e-QSL und EHam als englisches \ipa{i}.
Schrägstriche trennen alternative Formulierungen oder Listenelemente. Klammern geben den Zusatz der Quelle
wieder. Platzhalter \texttt{<\#\#\#\#\#>} beziehungsweise \texttt{<\#\#\#\#>} sind nicht mitzusprechen
und erscheinen in der Lautschrift als Auslassung.

\medskip
\textbf{LaTeX.} UTF-8-Datei; zweimal mit XeLaTeX kompilieren. Benötigt werden DejaVu Sans sowie die
üblichen LaTeX-Pakete der Präambel. Die IPA-Zeichen stehen direkt als Unicode im Quelltext;
es sind keine externen Bilder oder Datendateien erforderlich.
\end{document}
